\section{Zadanie 13}
\autor{Dominik Korsa}
Niech \(G\) będzie skończonym grafem z co najmniej jedną krawędzią.
Rozważmy łańcuch Markowa, gdzie przestrzeń stanów \(\Omega(G)\) jest rodziną
wszystkich zbiorów niezależnych w \(G\).
Krok ze stanu \(M_t\) do następnego stanu działa następująco.
Losujemy jednostajnie \(v \in V(G)\).
Jeśli \(v \in M_t\) to \(M_{t+1} = M_t \setminus \{v\}\).
Jeśli \(v \notin M_t\) to próbujemy dodać \(v\) do \(M_t\).
To znaczy, jeśli \(M_t \cup \{v\}\) jest zbiorem niezależnym w \(G\)
to \(M_{t+1} = M_t \cup \{v\}\), a w przeciwnym przypadku \(M_{t+1} = M_t\).
Udowodnij, że ten łańcuch jest nieprzywiedlny, nieokresowy i jego rozkład
stacjonarny jest jednostajny.

\subsection*{Rozwiązanie}
Niech \(n = |V(G)|\).

\subsubsection*{Nieprzywiedlność}
Aby udowodnić nieprzywiedlność tego łańcucha musimy udowodnić, że wszystkie
pary stanów (zbiorów niezależnych w \(G\)) się wzajemnie komunikują,
czyli dla każdej pary zbiorów niezależnych \(A, B \subseteq V(G)\)
istnieje takie \(k\), że \(P_{A,\,B}(k) > 0\).

Niech \(A = \{a_1, a_2, \dots, a_l\}\) i \(B = \{b_1, b_2, \dots, b_m\}\)
to pewnie zbiory niezależne. Następujący skończony ciąg kroków przechodzi ze
stanu \(A\) do \(\emptyset\), a następnie do \(B\):
\begin{itemize}
  \item Usuń \(a_1\) ze zbioru \(A\). Taki krok jest możliwy, bo zbiór
  niezależny po usunięciu wierzchołka pozostaje niezależny, i ma
  prawdopodobieństwo \(n^{-1}\).
  \item Podobnie usuwaj wierzchołki \(a_2\) do \(a_l\). Po tych krokach jesteśmy
  w stanie \(\emptyset\).
  \item Dodaj \(b_1\) do stanu \(\emptyset\). Ten krok ma prawdopodobieństwo
  \(n^{-1}\)
  \item Podobnie dodawaj wierzchołki \(b_2, \dots, b_m\). Każdy z takich kroków
  jest możliwy, bo wszystkie stany są podzbiorami zbioru niezależnego \(B\),
  więc również są zbiorami niezależnymi.
\end{itemize}
Zatem te \(k+l\) kroków, każde o prawdopodobieństwie \(\frac{1}{n}\) wystarcza,
aby przejść ze stanu \(A\) do \(B\), więc \(P_{A,\,B}(l+m) \geq n^{-(l+m)} > 0\).

\subsubsection*{Nieokresowość}
Niech \(A\) to maksymalny ze względu na inkluzję zbiór niezależny w \(G\).
Z treści zadania w \(G\) istnieje przynajmniej jedna krawędź, a skoro \(A\)
jest zbiorem niezależnym, to któryś z końców tej krawędzi, nazwijmy go \(v\),
nie należy do \(A\).

Będąc w stanie \(A\) wylosowanie wierzchołka \(v\) spowoduje pozostanie w
stanie \(A\) (z maksymalności \(A\)), więc: \(P_{A,\,A} \geq n^{-1} > 0\).

Definicja okresu stanu to
\(\operatorname{o}(v) = \gcd \mathcal{T}(v)\), gdzie
\(
  \mathcal{T}(v) = \{ t \in \mathbb{N}_1 \mid P_{v,\,v}(t) > 0 \}
\).\\
\(P_{A,\,A} = P_{A,\,A}(1) > 0\), więc \(1 \in \mathcal{T}(v)\), więc
\(\operatorname{o}(v) = 1\). Z lematu z wykładu:
\begin{quote}
  W nieprzywiedlnym łańcuchu Markowa wszystkie stany mają ten sam okres.
\end{quote}
W takim razie wszystkie stany mają okres \(1\), zatem łańcuch jest
\textbf{nieokresowy}.

\subsubsection*{Jednostajność rozkładu stacjonarnego}
Udowodnijmy, że macierz przejścia \(P\) łańcucha Markowa z zadania jest
symetryczna, czyli dla każdego \(A, B \in \Omega(G)\) zachodzi
\(P_{A,\,B} = P_{B,\,A}\).

\begin{itemize}
  \item Jeżeli \(A = B\) to \textbf{tr.}

  \item Jeżeli \(A\) różni się od \(B\) dokładnie jednym wierzchołkiem,
  powiedzmy \(B = A \cup \{v\}\), to jedynie wylosowanie \(v\)
  (z prawdopodobieństwem \(n^{-1}\)) przechodzi między tymi stanami.
  Skoro \(A\) i \(B\) są zbiorami niezależnymi, to ten krok się udaje.
  Zachodzi \(P_{A,\,B} = P_{B,\,A} = n^{-1}\).

  \item Jeżeli \(A\) różni się od \(B\) więcej niż jednym wierzchołkiem,
  to \(P_{A,\,B} = P_{B,\,A} = 0\), bo każdy krok dodaje lub usuwa co najwyżej
  jeden wierzchołek.
\end{itemize}

Z \hyperref[set07.task11]{zadania 11} wiemy, że skoro macierz przejścia
\(P\) jest symetryczna, to rozkład stacjonarny jest rozkładem jednostajnym.
